\section{Detailed source and protocol definitions}
\label{app:design}
\subsection{Timestamp-bounded adapters}
\paragraph{Borg.}
All 500 task-event shards of Google ClusterData2011-2 supply the compute
application \citep{google2011}. We replay pending/running states in timestamp
and source order, retaining unknown preceding age for time-zero records.
The target is first state change at 1/5/15/60 minutes, including failures,
evictions and kills as well as completion. Current job and known
running-machine memberships define the graph. The evaluated final prefix
contains 144.6 million events; 28 six-hour origins span elapsed days
22--28.75, with 3,760,001 task--origin queries and 105,077--171,603 active tasks
per origin. Those repeated queries are not independent replications.

\paragraph{BPI 2017.}
The application workflow \citep{vandongen2017} uses first
\texttt{A\_Pending}, \texttt{A\_Denied} or \texttt{A\_Cancelled} after creation
at 1/3/7/14 days. This is a documented application milestone, not repayment
or necessarily the final event. Last observed workflow handling supplies the
relation. We exclude retrospective trace attributes, eventual offer outcomes
and future assignments. The original final period comprises 61 daily noon
UTC origins during November--December 2016, with 114,051 case--origin queries
and 1,484--2,189 cases per cohort.

\paragraph{BPI 2012.}
A fresh workflow release \citep{vandongen2012} supplies an analogous target:
first \texttt{A\_APPROVED}, \texttt{A\_DECLINED} or \texttt{A\_CANCELLED}
after \texttt{A\_SUBMITTED}. We normalize these four activity names only to
reuse the replay adapter, retaining their original target definition. All
other activity names remain distinct; workflow lifecycle values are
normalized to lower case. Missing handlers share no entity-level factor.
This is additional temporal/source evidence, not a claim about another
independent bank. February 1--28, 2012 provides 28 final daily origins, with
23,749 case--origin queries, 793--889 cases per cohort, and
follow-up through March 14 00:00 UTC. Activity vocabulary is derived from
the individual fit prefix only.

In every source, observed events through the origin build features and
memberships; outcomes are attached separately. Entire longest-horizon windows
must be observable. Prefix state carries across date partitions without
revealing future events. Recorded event time is the availability proxy;
the archives do not establish live ingestion latency or exclusive resource
ownership. The municipal BPI 2015 candidate was deferred after a
development-source audit did not establish a sufficiently reliable closure
definition across its activity labels. An Alibaba 2018 task-table candidate
was also deferred: absent end timestamps dominated some reconstructed active
cohorts, and their interpretation as continued running was unverified.
Neither candidate contributed predictive results.

\subsection{Initial frozen evaluation}
The initial study compares HGB personal/global/relational arms and temporal
Mixer, GRU and tabular MLP ensembles. Borg uses 32 base features plus a
144-event-count extension; BPI 2017 uses 103 features. Temporal models use
the last 16 events per task/job/machine or case/handler scope, width 64 and
a five-category first-exit head. Selection uses equal-origin mean Brier over
the four horizons, including monotone calibration variants and three-seed
neural means; no best-seed selection is permitted.

Borg individual training/calibration/development end at days 18/20/22;
dependence uses forward-fold predictions in days 8--10, 12--14 and 16--18,
each produced by earlier fitted and calibrated predictors. BPI 2017 uses
training before September, September calibration and October development,
with June/July/August forward folds for dependence. Selected individual
models are HGB global for pending tasks, HGB relational for running tasks,
and the temperature-calibrated Mixer ensemble for BPI 2017. Gaussian typed
and global models are selected for Borg and BPI respectively.

An immutable 264-file manifest froze code, models and evaluation choices
before those final periods were opened. This was a development-informed
holdout freeze, not external preregistration. A lossless Borg timestamp-type
repair after a replay failure retained all original data and changed no
timestamp value or frozen model. Original scientific results and frozen assets
remain preserved; obsolete manuscript copies are not part of the current paper.

\subsection{Retrospective diagnostics and fresh confirmation}
The initial outcomes motivated the direct DM controls, delayed calibration,
factor removals and joint-profile scores. Their analyses on the original
final periods are therefore \emph{retrospective}. Direct count training uses
24 complete Borg six-hour OOF cohorts and 50 complete BPI 2017 OOF cohorts.
The stronger development-fitted DM control instead uses the eight Borg or
17 BPI 2017 current-version development origins, matching the historical
residual fitting period. Both fitting regimes are retained.
The original residual comparator fits on development residuals,
with leave-one-origin-out development scoring and three sampling seeds.

For fresh confirmation, BPI 2012 individual fitting ends before November 1,
2011; November calibrates the individual model, December fits dependence and
direct counts, and January selects configurations. The global-feature HGB
with identity correction wins individual Brier. Identity/global wins
Gaussian development count CRPS; bias-DM wins the direct comparison;
bias-DM conditioning wins the conditioning comparison. Static and rolling
correction remain prespecified controls, with dependence fitted separately
for each treatment. All six Gaussian fits per seed use only development
data.

Before retaining any February outcomes, a second immutable manifest froze
82 assets, the primary comparisons, metrics, origins and bootstrap settings.
The raw parser had previously discarded every event at or after February 1
before retaining its activity or outcome. Final predictors remain fixed;
only the explicitly prespecified rolling correction may use earlier
February outcomes after their 14-day maturity delay. January 19--31 origins
provide cross-boundary calibration history but are not final score origins.
We retain all prespecified outcomes irrespective of which method wins.
\subsection{Cross-domain extension and stronger graph controls}
We next screened additional public sources with bounded local storage, retaining
unsuccessful and unsuitable candidates. Azure V2 \citep{azure2019} supplies a
second compute provider. A first 49,231-row schema/prediction pilot was followed
by a sample retaining \emph{all released VMs} in subscriptions selected by a
fixed SHA-256 rule (first 32 bits modulo 20 equal zero). The sample is complete
within those released subscriptions, not the provider's full production
population. The target is VM deletion at 5/30/120/360 minutes; observed deployment
and subscription memberships supply two relations. Lifetime CPU summaries,
usage-derived VM categories, and deletion times are excluded from features.
Time-zero creations give an observed age lower bound; timestamp coarsening and
the publisher's workload sampling limit interpretation. The interval adapter
emits creation events only, so Azure temporal tokens summarize prior creations,
not a full stream of VM deletions or CPU measurements.

The UCI incident log \citep{amaral2018} supplies handler and team relations.
We predict the first \emph{recorded} Resolved/Closed update after the first
observed nonterminal update at 1/6/24/72 hours. This is an event-log milestone,
not an estimate of an unobserved physical repair time. Resolution timestamps,
closure codes, eventual resolvers, and SLA outcomes repeated in earlier rows
are excluded. BPI 2014 \citep{vandongen2014activity} supplies timestamped incident
activities and team membership; its target is the first Closed activity after
an observed Open. We do not reconstruct live arrival populations from either
historical case extract. Event timestamps proxy information availability.
Naive incident timestamps are interpreted on a UTC clock for reproducible
partitioning; the archives do not verify operational timezone or ingestion
latency. This convention does not recover daylight-saving transitions.

All sources have distinct individual-fit, calibration, dependence-fit, and
development-validation windows. Entire longest-horizon labels must mature
strictly before each window ends. Azure uses day boundaries 0/10/13/16/20;
UCI uses March 1/22, April 1/15, and May 1, 2016; BPI 2014 uses October 1,
November 1/15, December 1, 2013, and January 1, 2014. Prefix state carries
across boundaries. New-period confirmation fixes 28 origins per source:
Azure days 22--28.75, UCI May 1--28, and BPI 2014 January 1--28.
An executable asset manifest freezes all fitted models and comparisons before
these targets are prepared or scored. Earlier schema audits inspected whole-file
state totals and time ranges; this is an internal prediction-holdout freeze,
not external preregistration or a claim that the raw files were never opened.

\paragraph{Matched individual models.}
Controls include personal-feature HGB, HGB with current relation summaries,
and HGB with the same observed event-history tokens available to temporal
attention. Neural controls use width 48: a node MLP, a global set model,
two rounds of typed entity--group mean message passing, and temporal attention
over the last eight observed relation events followed by message passing.
These are independent survival-target adaptations, not reproductions of TGN,
TGAT, or DyGFormer benchmark systems \citep{rossi2020,xu2020,yu2023}.
Each network uses a five-category first-exit head. UCI trains for 30 epochs;
Azure and BPI 2014 for 20, with AdamW learning rate .001 and weight decay .001.
HGB uses 100 boosting iterations, 15 leaves and L2 penalty 10. Checkpoints and temperature
are chosen on the earlier calibration window; three-seed mean CDFs determine
the selected architecture. Relation-specific and size-preserving shuffled
graph controls keep the input feature set and training budget fixed.
A separate bounded Borg pilot shares the existing 176 engineered features
across models, testing additional message passing beyond those summaries;
its sampled-node scores are not population estimates.

\paragraph{Joint controls.}
Gaussian factors include independent, global, each relation alone,
global-plus-each-relation, combined typed, and shuffled-incidence controls.
Small development-only variance grids and coordinate search fit count scores;
the typed model may reduce to the global control. Coordinate fitting starts
at the global fit and performs two sweeps over weights; it is not an exhaustive
joint-grid optimum. Separately fitted relation-only arms are retained even
when they beat that combined fit. Student-$t$ copulas add a
shared inverse-chi-square scale and compare global and typed variants with
degrees of freedom chosen from 4/8/16. Each draw still uses one entity uniform
across horizons. Concentration-only and bias-adjusted direct DM, exact
beta-binomial count scores, and historical residual forecasts use the same
dependence-fitting period. None is tuned on confirmation.

A small additional comparator uses the unchanged published TACTiS attentional
copula module \citep{drouin2022}, with a custom current-cohort encoder,
randomized discrete probability-integral-transform targets, and rank-adjusted
generated uniforms. All comparators in that pilot use the same 64-node subsets.
It is explicitly a limited adaptation: neither the full time-series model nor
an exact discrete likelihood reproduction, and not evidence of superiority
over TACTiS or TACTiS-2 in their original tasks.

Exploratory fitting uses empirical-ensemble CRPS; confirmation uses the fair
distinct-pair estimator with 512 draws and three sampling seeds. The peak
reference scale in confirmation is $\sqrt{1+\sum_iF_i(1-F_i)}$, as in the
original study. Development pilot peak scores used a max-with-one stabilization
and are not numerically pooled with confirmation. Seeds are averaged within
origin before 2,000 paired circular block replicates with 1/3/7-day blocks.
The same frozen probabilities and reference risk quantities are used for
every joint comparator.

\paragraph{Screened alternatives and limits.}
A Google 2019 \citep{google2019} event-shard audit found approximately 83\% of
sampled entity histories contained one event; that shard cannot support a
complete-trajectory forecast. BPI 2014 change snapshots lack version histories
for planned times and affected configuration items; incident activities offer
a better-supported relation replay. Citi Bike \citep{citibike} and airline
\citep{bts2024} pilots test transportation at small scale. Published trips are
retrospectively filtered, and airline schedule-version availability is an
explicit assumption. Their exploratory scores are not presented as live
deployment results. Source selection followed development feasibility and
performance, so reporting all screened candidates is part of the evidence.

